Vision-Language Models (VLMs) have achieved remarkable performance across multimodal tasks, yet their robustness against adversarial perturbations remains poorly understood. This paper presents the first comprehensive benchmark evaluating adversarial robustness across 16 state-of-the-art VLMs spanning six model families (Qwen, SmolVLM, InternVL, LLaVA, DeepSeek, and others) with parameters ranging from 256M to 8.3B. We systematically evaluate 9 adversarial attack types: 5 white-box methods (FGSM, PGD, AutoPGD, AutoConjugateGradient, BasicIterative) and 4 black-box methods (Square, SimBA, Boundary, SignOPT) across 3 epsilon levels (0.02, 0.05, 0.08) on 8 diverse vision-language tasks. Our novel epsilon-based attack framework provides direct perturbation control without optimization, enabling fast and reproducible adversarial evaluation.

Our findings reveal significant vulnerabilities across all model families, with accuracy degradation ranging from 15.2\% to 47.8\% under adversarial perturbations. Smaller models exhibit higher variance in robustness, while larger models demonstrate more consistent but still substantial vulnerabilities. We identify architecture-specific robustness patterns: transformer-based models show better resilience to black-box attacks, while CNN-hybrid architectures are more susceptible to white-box methods. The comprehensive evaluation across 8 diverse vision-language tasks (chart interpretation, table extraction, spatial reasoning) reveals task-dependent vulnerability patterns. Our analysis of memory-performance trade-offs provides practical insights for deploying VLMs in resource-constrained environments. This benchmark establishes the first systematic framework for VLM robustness evaluation and highlights the urgent need for robust multimodal architectures.